\documentclass{umk-matematika}
\begin{document}

\begin{center}
{\Large\bfseries Практическая работа 05}
\end{center}
\vspace{0.5cm}

\textbf{Задача 1.} Сформулируйте аксиому 1 стереометрии.\par\vspace{10pt}
\textbf{Задача 2.} Сформулируйте аксиому 2 стереометрии.\par\vspace{10pt}
\textbf{Задача 3.} Сформулируйте аксиому 3 стереометрии.\par\vspace{10pt}
\textbf{Задача 4.} Сколько плоскостей можно провести через две пересекающиеся прямые?\par\vspace{10pt}
\textbf{Задача 5.} Могут ли три точки одновременно принадлежать двум разным плоскостям?\par\vspace{10pt}
\textbf{Задача 6.} Даны четыре точки, не лежащие в одной плоскости. Сколько различных плоскостей можно провести через различные тройки этих точек?\par\vspace{10pt}
\textbf{Задача 7.} Что такое сечение многогранника? Какие многоугольники могут быть сечениями куба?\par\vspace{10pt}
\textbf{Задача 8.} Постройте сечение куба $ABCDA\_1B\_1C\_1D\_1$ плоскостью, проходящей через середины рёбер $AB$, $BC$ и вершину $D\_1$. Какой многоугольник получится?\par\vspace{10pt}
\textbf{Задача 9.} Какие многоугольники могут быть сечениями тетраэдра?\par\vspace{10pt}
\textbf{Задача 10.} Строительный блок имеет форму прямоугольного параллелепипеда. Его распилили плоскостью, проходящей через диагонали двух противоположных граней. Какая фигура получилась в сечении?\par\vspace{10pt}
\newpage
\section*{Ответы}

\begin{enumerate}
  \item Ответ: аксиома 1.
  \item Ответ: аксиома 2.
  \item Ответ: аксиома 3.
  \item Ответ: одну.
  \item Ответ: да, если точки лежат на одной прямой; нет, если не лежат.
  \item Ответ: 4 плоскости.
  \item Ответ: треугольник, четырёхугольник, пятиугольник, шестиугольник.
  \item Ответ: пятиугольник.
  \item Ответ: треугольник и четырёхугольник.
  \item Ответ: прямоугольник.
\end{enumerate}
\end{document}
